\section{Introduction}

For a chain complex of free abelian groups,
\[
  \cdots \longrightarrow C_{n+1}
  \xrightarrow{d_{n+1}} C_n
  \xrightarrow{d_n} C_{n-1}\longrightarrow\cdots,
  \qquad d_nd_{n+1}=0,
\]
ordinary homology is
\[
  H_n=\ker d_n/\operatorname{im}d_{n+1}.
\]
Smith normal form supplies the standard finite-chain algorithm for computing
integer homology \citep[Chapter 3]{kaczynskiMischaikowMrozek2004}, but computing
$H_n$ requires more than applying Smith form independently to the two
boundaries.  One must extract an integral basis of the kernel, express the
incoming boundary in exactly that basis, prove that the resulting
coordinates are valid, and relate the second cokernel to the original
quotient.

These steps are especially important in a proof assistant.  A rank count
alone cannot supply the change of basis.  A determinant-based invertibility
witness does not expose the inverse needed to transport coordinates.  A
runtime list of invariant factors is not itself a theorem about the quotient.
Index arithmetic also has to cover zero groups and the top and bottom
degrees without informal dimension conventions.

We implement the complete construction in Lean 4~\citep{deMouraUllrich2021Lean4}
and mathlib~\citep{mathlib2020}, using the strong Smith-normal-form results
of \code{lean-normal-forms}.  Each Smith result records $U$, $V$, chosen
inverses, both inverse identities, and
\[
  Ud_nV=S.
\]
The main contributions are:

\begin{enumerate}
\item a finite-free integer chain-complex interface with a fixed matrix
  orientation and an intrinsic homology quotient;
\item an explicit linear equivalence between Smith tail coordinates and
  $\ker d_n$;
\item a proof that $V^{-1}d_{n+1}$ has zero non-kernel prefix;
\item a second certified Smith calculation and an additive decomposition of
  ordinary homology into cyclic and free summands;
\item a finite normalized simplicial frontend deriving signed boundary
  matrices from nondegenerate face tables;
\item an exact comparison with mathlib categorical homology, including the
  exceptional shape of the degree-zero short complex;
\item chain isomorphisms with mathlib's normalized nondegenerate and
  normalized Moore complexes from an explicit simplicial-realization
  witness;
\item a theorem-connected executable projection, independent facade, native
  regression, axiom audit, and reproducible artifact.
\end{enumerate}

The v1.2.2 scope covers explicit boundary matrices, finite normalized
nondegenerate-simplex tables, and their comparison with mathlib's abstract
homological-complex infrastructure.  The normalized face table stores less
structure than a full simplicial set; accordingly, the normalized Moore
result assumes a checked realization witness rather than manufacturing
missing simplicial identities or degeneracy maps.
